\documentclass[12pt,onecolumn]{article}

\usepackage[utf8]{inputenc}
\usepackage[T1]{fontenc}
\usepackage{lmodern}
\usepackage[margin=1in]{geometry}
\usepackage{setspace}
\usepackage{microtype}
\usepackage{amsmath,amssymb}
\usepackage[numbers,sort&compress]{natbib}
\usepackage[colorlinks=true,linkcolor=blue,citecolor=blue,urlcolor=blue]{hyperref}
\usepackage{orcidlink}
\usepackage{fancyhdr}

\newcommand{\doi}[1]{\href{https://doi.org/#1}{#1}}
\setstretch{1.10}
\pagestyle{fancy}
\fancyhf{}
\setlength{\headheight}{15pt}
\lhead{Motion Is Relational}
\rhead{John C. W. McKinley}
\cfoot{\thepage}

\title{
\textbf{Motion Is Relational}\\
\large A No-Go on Independent Position and Substance Displacement\\[0.5em]
\large The Universal Present, Local Accrual, and Lawful Physical Relations
}
\author{John C. W. McKinley\,\orcidlink{0009-0005-7097-5035}}
\date{October 2026}

\begin{document}
\maketitle

\begin{abstract}
Motion is conventionally described as the displacement of a physical object through space. Although this description is mathematically successful, its ordinary interpretation invites an additional ontology: independently existing objects occupying successive positions within an independently existing spatial container. This paper argues that none of these independent entities is required. Building on the relational account of physical actuality developed in \emph{Space Was Never a Completed Container}, it proposes that physical reality consists of lawful configurations and relationships, while what are ordinarily called objects are stable, identifiable patterns of lawful organization.

Under this account, position is not an intrinsic property of an object but a relational assignment. Bob is stationary in Bob's own coordinates, and Alice is stationary in Alice's own coordinates. When Bob and Alice are described as moving relative to one another, neither possesses an absolute velocity or undergoes displacement through an independently existing spatial background. What changes is their lawful spatial relationship. That change is physically real, and its consequences are registered through interactions, including changing separation, Doppler shifts, acceleration measurements, and eventual physical encounters. The measurements do not create the relationships; they register aspects of an evolving lawful configuration.

The same relational account extends to time and physical actuality. There is one universal now: the universe has one current physical configuration, not an independently existing sequence of past, present, and future configurations. This universal present is an ontological statement about physical actuality, not the assertion of a universal clock or privileged coordinate system. Proper time is local to each timelike worldline, with its accumulated amount determined by the worldline and the governing spacetime metric. Different physical constituents therefore possess different accrued proper-time histories while participating in the same physical actuality. The invariant information speed limit restricts the propagation of information and what observers know about distant conditions; it does not determine whether those distant conditions presently exist.

The paper distinguishes invariant quantities associated with individual worldlines, including proper acceleration and accumulated proper time, from spatial and temporal relationships assigned between physical systems. It advances a conditional no-go: given the absence of independently existing spatial positions and traveler-substances, independent substance displacement has no physical referent. Motion remains a valid description of lawful relational change, but it does not establish the existence of a substance traveling through a completed spatial container.

The result changes no equations or predictions of special or general relativity. Its contribution is ontological: there is one present physical actuality, constituted by lawful relationships whose configurations change, while individual constituents accumulate their own proper-time histories. Neither motion nor time requires an independently existing spatial or temporal container.
\end{abstract}

\section{Introduction}

The statement that an object moves from one location to another is an ordinary and useful description of physical events. Astronomers calculate planetary trajectories, engineers track spacecraft, and observers measure the changing separation between physical systems. These descriptions are successful because the mathematical relationships used to predict physical interactions are successful. The present paper does not dispute those calculations or their observable consequences.

The question addressed here is what those calculations require us to believe exists.

The conventional mental picture places objects inside a spatial container. The container has locations, the objects occupy those locations, and motion consists of objects passing from one location to another. Time provides the corresponding temporal container within which these movements occur. Even when modern relativistic mathematics replaces Newtonian absolute space and time, the older mental picture frequently survives in ordinary interpretation.

The physical object is imagined to remain an independently existing substance while changing its position within an independently existing spatial arrangement. The trajectory is then treated as the path through which that substance travels.

This paper questions the necessity of that picture.

The relational account developed in \emph{Space Was Never a Completed Container} \cite{SpaceContainer} holds that space is not an independently completed physical container awaiting occupation by objects. What exists physically is the lawful relational configuration of physical systems and the interactions through which that configuration is expressed. Spatial descriptions organize those relationships, but they do not establish an additional spatial substance underlying them.

Once that premise is adopted, the conventional interpretation of motion must be reconsidered. If independently existing spatial positions are absent, an object cannot possess an absolute location. If what are ordinarily called objects are stable patterns of lawful organization rather than independent traveler-substances, the movement of a substance through an independent spatial container is no longer required to account for observed physical change.

What remains is lawful relational change.

The same reasoning extends naturally to temporal description. Proper time is accumulated locally along individual timelike worldlines, while different worldlines acquire different accumulated intervals. There is no requirement that all physical constituents agree about how much proper time has elapsed. Nor does the absence of a universal elapsed-time quantity establish the existence of multiple physically separate presents.

The present paper therefore develops a unified interpretation of position, motion, proper-time accrual, and physical actuality. Its central proposal is that there is one current physical actuality, consisting of lawful relationships and configurations, without an independently existing spatial container, a universal master clock, or a physically retained past.

The result is deliberately interpretive. The paper does not introduce a new field equation, modify relativistic predictions, or claim that the relational ontology is experimentally established by the familiar absence of absolute velocity. It instead examines the physical description that remains when independent spatial position and substance displacement are removed from the ontology.

\section{Physical Actuality Without an Independent Spatial Container}

The conventional representation of a physical system ordinarily begins by assigning it a position. The position is then expressed relative to an origin and coordinate axes. Other systems receive corresponding positions, and their separations and movements are calculated from the resulting mathematical relationships.

Nothing in this procedure requires the coordinate system to be an independently existing physical object.

A spatial coordinate grid is a map. Its origin, orientation, and coordinate assignments are selected for the purpose of describing relationships among physical systems. Another coordinate grid may be selected without changing the physical situation.

For example, an astronomer may choose one star as the origin and assign positions to twenty neighboring stars. The resulting arrangement is useful and predictive. A second astronomer may choose a different star as the origin and construct a different map. The coordinate values change, while the physical relationships being represented remain the same.

The map does not establish a master grid existing independently of the stars.

This distinction is already familiar in relativity, where physical descriptions must remain meaningful under changes of coordinate representation \cite{Wald,Einstein1905}. The present paper extends the distinction into an explicit ontology.

Under the relational interpretation, what is physically actual is not a collection of objects occupying independently existing positions. Physical actuality consists of lawful relationships, interactions, and configurations. Position identifies a relationship within a specified description. Distance identifies a metric relationship. Direction and velocity likewise require an appropriate relational standpoint.

The absence of an independently existing spatial container does not imply the absence of spatial measurement. On the contrary, spatial measurement remains an indispensable means of describing the lawful relationships that constitute physical actuality.

A ruler measures the relationship between its own physical structure and another physical system. Radar establishes relationships through the exchange of signals. A telescope registers radiation received from distant sources. These procedures produce meaningful measurements because the governing relationships are lawful and reproducible.

Their success does not require an additional spatial substance in which the measured systems reside.

The governing laws specify the admissible relationships, interactions, and outcomes. What is ordinarily called a physical configuration is the relational condition expressed by physical actuality, not an additional substance underlying it.

\section{Objects as Stable Lawful Patterns}

The relational interpretation also changes the meaning assigned to a physical object.

An object is ordinarily imagined as a persistent substance that possesses physical properties and occupies spatial locations. Its continued existence is represented as the continued presence of that substance while its location or other properties change.

The present model instead identifies an object as a stable, recognizable organization of lawful physical relationships.

An atom, for example, is identified by its lawful structure and characteristic responses to physical interactions. Its observable properties are established through the conditions governing those interactions. The matter-side account developed in \emph{The Atom Is Available When Called On} \cite{Atom} describes a bound system as a lawful stationary structure whose physical consequences are registered under admissible interactions.

The present paper adopts that distinction between lawful organization and an independently traveling substance.

Calling an atom an object remains convenient. Calling a planet an object remains convenient. Calling Bob or Alice an object remains convenient. The terms identify organized physical systems whose lawful characteristics support recognition and prediction.

What the relational model removes is the additional requirement that each identified system consist of an independent traveler-substance occupying successive positions within a completed spatial container.

Stable organization does not mean an absence of physical interactions or changes. A living body, for example, contains continuously interacting constituents. Its persistence as an identifiable system consists of the continued lawful organization of those constituents, not the absence of internal physical activity.

The same distinction applies to stars, planets, molecules, and other identifiable physical systems.

A physical system is therefore not made unreal by removing the picture of an independently existing substance. Its lawful properties, interactions, and continuing organization remain physically consequential.

The model rejects independent substance ontology, not physical actuality.

\section{Position Is Relational and Motion Is Relational Change}

Consider two physical systems, Bob and Alice, initially separated from one another.

Suppose Bob and Alice approach each other.

In a coordinate system comoving with Bob, Bob remains at the spatial origin while Alice's assigned position changes. In a coordinate system comoving with Alice, Alice remains at the spatial origin while Bob's assigned position changes.

Neither observer possesses an absolute spatial velocity.

The same physical relationship admits both descriptions. Each observer is stationary in the coordinates constructed around that observer's own worldline, while the other observer receives changing spatial coordinates.

This is a standard feature of relativistic kinematics, not a new physical postulate \cite{TaylorWheeler,Einstein1905}.

For accelerated observers, the comoving description is understood locally or through coordinates adapted to the observer's worldline. A single inertial frame need not remain comoving throughout an accelerated trajectory. This does not affect the present point: an observer's own worldline may be represented at fixed spatial coordinates without eliminating the measurable relationships between that observer and other systems.

The relational ontology asks what, beyond these descriptions, is physically required.

Bob registers Alice as approaching. Alice registers Bob as approaching. Their separation changes, their exchanged signals exhibit predictable relationships, and they may eventually meet.

These physical consequences establish that the relationship changes. They do not establish that either Bob or Alice possesses an absolute position from which an independent substance has departed.

The change belongs to the relationship.

This distinction is not merely verbal. The proposition that Bob changes his absolute position requires an independently specified positional background. Without such a background, the proposition lacks the additional physical referent it purports to describe.

The statement that Bob's relationship to Alice changes requires no such background. It refers directly to the lawful physical relationship between identifiable systems.

Nor does the relational interpretation reduce motion to an illusion or to an observer's subjective report.

The relationship itself changes. Instruments register that change. The instrument readings are physical consequences of the lawful configuration; they do not create the configuration by observing it.

The motion description is therefore real and operationally indispensable, while independent spatial displacement is unnecessary.

The distinction may be expressed simply: neither Bob nor Alice undergoes intrinsic spatial displacement through an absolute background. What evolves is their lawful relationship, and that evolution is what the language of relative motion describes.

\section{A No-Go on Independent Substance Displacement}

The preceding account supports a conditional no-go.

Independent substance displacement requires an independently existing traveler-substance, independently existing spatial positions, and a physically meaningful passage of the traveler between those positions.

The relational ontology supplies none of these independent entities.

What it supplies is lawful physical organization, spatial relationships among physical systems, and lawful changes in those relationships.

Consequently, independent substance displacement has no additional physical referent within the model.

This conclusion is not obtained merely by defining a coordinate system in which an observer is stationary. Such a definition would establish only the familiar fact that an observer occupies fixed spatial coordinates in a comoving description.

The present conclusion depends instead on the prior ontological premise that spatial positions and traveler-substances do not exist independently of the lawful physical relationships being described.

Once that premise is adopted, the conventional traveler interpretation introduces entities not required by the physical account.

An observer may describe a spacecraft as crossing a million kilometers. Radar measurements may establish the spacecraft's changing range. The spacecraft may arrive at a destination and produce a collision or docking event. All of these are real physical consequences.

What the model denies is that those consequences require a substance to have traversed an independently existing spatial container.

The spatial interval is a lawful metric relationship. Its evolution describes the changing physical configuration. The endpoint interaction is another lawful physical event.

The additional traveling-substance picture contributes no necessary physical quantity to the calculation.

Thus the no-go is ontological and conditional. It follows from the relational account of physical actuality, not from the equations of relativity alone.

Its significance lies in eliminating an unnecessary interpretation while retaining the entirety of the measurable physical behavior.

\section{Acceleration, Worldline Geometry, and Real Physical Change}

The absence of independent substance displacement does not eliminate acceleration.

An accelerating physical system registers a genuine local physical quantity. An accelerometer carried by a rocket detects the rocket's proper acceleration without requiring comparison to a distant observer.

In relativistic notation, the four-acceleration is

\begin{equation}
a^\mu=\frac{D u^\mu}{d\tau},
\end{equation}

where \(u^\mu\) is the four-velocity and \(\tau\) is the proper time along the worldline.

The magnitude of proper acceleration is an invariant physical quantity.

This separates acceleration from ordinary spatial velocity.

Spatial velocity is assigned relative to a reference frame. Proper acceleration is measurable locally along the accelerated worldline.

When Alice fires her rocket, the lawful conditions governing her subsequent relationships with other physical systems change. Her worldline acquires a different geometric development from the worldline she would have followed without the rocket burn.

Her accelerometer records a real physical effect. Her future separation from Bob develops differently. Their subsequent exchange of signals reflects that difference.

None of this requires Alice to change her spatial location within her own adapted coordinate description.

The distinction is particularly clear in gravitational situations.

A person standing on the surface of a planet experiences nonzero proper acceleration because the supporting surface prevents free fall. A freely falling observer experiences zero proper acceleration, neglecting nongravitational forces, while following a geodesic of the surrounding spacetime geometry.

The supported observer and the freely falling observer follow different worldlines. Their instruments register different local conditions, and their accumulated proper times depend on the metric along their respective paths.

Proper acceleration is therefore a genuine feature of physical behavior, while independent spatial displacement remains unnecessary.

The relational model preserves acceleration as a lawful local fact and relative motion as lawful relational change.

\section{Proper Time as Local Accrual}

The same distinction between intrinsic and relational quantities applies to elapsed time.

A massive physical system follows a timelike worldline and accumulates proper time along that worldline. The accumulated interval is determined by the spacetime metric and the worldline's geometric development.

Using the metric signature \((-+++)\), proper time satisfies

\begin{equation}
c^2d\tau^2=-g_{\mu\nu}dx^\mu dx^\nu.
\end{equation}

In flat spacetime, relative to a specified inertial coordinate system, this becomes

\begin{equation}
d\tau=dt\sqrt{1-\frac{v^2}{c^2}},
\end{equation}

where \(v\) is the system's spatial velocity in that coordinate system and \(c\) is the invariant speed of light.

The equation shows that local proper-time accrual and relative spatial velocity are related through the spacetime metric.

At zero relative velocity, the proper-time interval equals the coordinate-time interval in the selected inertial system. As relative velocity increases, the proper-time interval per unit coordinate time decreases.

This does not mean that the moving observer detects an abnormal rate of operation in the observer's own clock. A local clock proceeds according to its own proper time.

The distinction becomes apparent when different worldlines are compared.

Suppose Bob and Alice depart together from a space station in deep space. Bob remains with the station while Alice travels outward, reverses direction, and returns.

If Alice spends most of her journey traveling at a speed of \(0.8c\) relative to the station, the proper-time accrual during those coasting portions is

\begin{equation}
d\tau_A=0.6\,dt_B.
\end{equation}

In the idealized case in which the changes of velocity occupy negligible intervals, ten years on Bob's clock correspond to six years on Alice's clock.

The difference is accumulated along their distinct worldlines. It does not imply that Alice's clock malfunctioned or that she temporarily ceased participating in physical actuality.

The twin result is a relationship between worldline histories.

The locally accrued proper time of each participant remains the interval measured along that participant's own worldline, consistent with the local-accrual interpretation developed in earlier work \cite{Accrual,DilationProperty,WorldlinesProduce}.

The difference in accumulated proper time does not establish a master clock against which all physical systems must be judged.

It establishes that distinct worldlines connecting common events possess different proper-time lengths.

\section{The Motion Tax and Gravitational Conditions}

For explanatory purposes, the difference between a selected reference-time interval and the proper time accumulated along another worldline may be called a relativistic motion tax.

The term does not identify a new force or physical mechanism. It is an accounting description of the difference between a reference interval and the locally accumulated interval.

In the idealized space-station example, Alice's travel produces a motion tax relative to Bob's reference clock.

Gravity introduces a further contribution.

In a weak, approximately stationary gravitational field, the relationship between proper time and a selected coordinate time is approximately

\begin{equation}
\frac{d\tau}{dt}\simeq
1+\frac{\Phi}{c^2}-\frac{v^2}{2c^2},
\end{equation}

where \(\Phi\) is gravitational potential and \(v\) is coordinate velocity in the chosen weak-field description.

The equation displays both the gravitational contribution and the contribution associated with relative motion.

A clock deeper in a gravitational potential well generally accumulates less proper time than a suitably compared stationary clock higher in that well. A clock moving relative to the selected reference coordinates also acquires a velocity-dependent adjustment.

In the space-station example, Alice may depart from Bob, travel to a planet, remain on its surface, and eventually return.

Her accumulated proper time reflects the entire worldline, including her high-speed travel and the interval spent at the planet's gravitational potential.

Standing on the planet means resisting free fall. That resistance produces proper acceleration. The gravitational contribution to clock accrual, however, is not identical to that proper acceleration. A freely falling clock is also governed by the gravitational metric.

The important unifying point is that the spacetime metric determines local proper-time accrual along each worldline.

The invariant speed \(c\) enters both the velocity-dependent relationship and the gravitational relationship. This motivates an interpretation in which the causal structure expressed by \(c\) is fundamental to understanding local accrual.

The value of \(c\) alone does not determine every proper-time interval; the metric and the worldline supply the remaining geometric information.

Thus proper time is not treated as a universal substance that passes independently of physical conditions. It is a locally accumulated invariant quantity associated with lawful physical history.

The motion-tax description remains useful because it illustrates the difference between a chosen reference interval and the interval accumulated along a particular worldline, without requiring a universal master clock.

\section{One Universal Present}

The relational interpretation of space and motion invites a corresponding examination of physical time.

The present paper adopts the proposition that there is only one physical actuality, which we call now.

The universe does not consist of an independently existing past, an independently existing present, and an independently existing future, each retained as a completed physical configuration. There is one current physical configuration, constituted by lawful relationships and physical organization.

The past does not remain physically in existence as another universe.

Records, memories, photographs, fossils, radiation, and other evidence of earlier events exist within current physical actuality. The events to which those records refer are not thereby retained as separately existing physical occurrences.

A photograph of an earlier event is a presently existing physical record. The photograph does not preserve the earlier event as an independently continuing physical configuration.

Likewise, a human memory is a present physical condition of the remembering system. Its content concerns earlier events, but the memory itself exists now.

The physical universe therefore contains present records from which earlier physical configurations may be reconstructed. It does not follow that those earlier configurations continue to exist alongside the current one.

This conception is consistent with the broader argument developed in \emph{The Universe Has No Past} \cite{NoPast}, which distinguishes the physically actual configuration from reconstructed histories of earlier configurations.

The universal present proposed here is not an independently existing temporal container. It is the name given to physical actuality itself.

The model does not require a universal quantity of elapsed time to govern that actuality. It does not require all physical systems to have accumulated the same proper time. Nor does it require a completed spatial grid on which the current configuration must be placed.

There is one actuality, expressed through lawful physical relationships.

The absence of a universal clock does not entail the absence of universal actuality.

\section{Local Proper Time Does Not Divide the Present}

Consider again Bob and Alice.

They depart together, follow different worldlines, and reunite. Bob's reference clock has accumulated thirty-eight years, while Alice's reference clock has accumulated thirty-one.

Their accumulated proper times differ by seven years.

Nevertheless, Bob and Alice meet and shake hands.

The physical interactions involved in that handshake proceed according to the local lawful states and relationships of the participating systems. The interaction does not require the two participants to possess equal accrued proper-time histories.

Bob does not have to surrender seven years of history before the interaction is permitted. Alice does not have to acquire seven years of additional history. The laws governing the encounter do not require a universal calendar reconciliation.

Their different histories participate in one physical interaction.

The example becomes more revealing when the physical constituents of their bodies are considered individually.

Bob does not consist of a single worldline. His body contains many constituent systems with distinct physical histories. The same is true of Alice.

Individual cells, and the material constituents composing them, follow their own timelike histories. Their accumulated proper times need not be exactly equal, even within one body.

For purposes of ordinary description, the statement that Bob is thirty-eight years old identifies a representative history, commonly associated with elapsed time measured by a suitable reference clock. It does not establish that every physical constituent of Bob's body has accumulated precisely the same proper-time interval.

Yet these constituents participate in a coherent living physical organization.

Different accumulated histories do not prevent lawful interaction within Bob's body. They do not prevent interaction within Alice's body. And they do not prevent Bob's and Alice's bodies from interacting when they meet.

The result is not that every constituent literally occupies the same spacetime event. Rather, local interactions among constituents collectively establish the organized physical relationships described as bodies, encounters, and handshakes.

The amount of proper time accumulated by one constituent is not a prerequisite that must equal the amount accumulated by another.

Local proper-time accrual identifies the history of a worldline. It does not partition physical actuality into separately existing presents.

Under the ontology proposed here, Bob and Alice participate in the same current actuality despite carrying different accumulated histories.

The same statement applies to the constituents of their bodies and to other physical systems throughout the universe.

There is one present physical actuality, with many local histories.

\section{The Information Speed Limit and Distant Actuality}

The invariant causal speed limit places a fundamental restriction on the propagation of information between physical systems.

A distant physical condition cannot ordinarily be reported instantaneously to a remote observer. Electromagnetic signals require propagation intervals, and the information available to an observer is limited by the observer's causal relationships with other events.

This limitation is central to relativistic physics.

It does not, by itself, establish that distant physical actuality is absent.

Suppose Bob observes a distant star through a telescope. The light reaching Bob records conditions associated with the star at the emission events responsible for that light. Those conditions are not necessarily the conditions Bob would describe as current.

Bob therefore lacks immediate information about the distant star.

The absence of that information is not evidence that the star lacks a current physical condition.

The distinction is between physical actuality and the availability of information about physical actuality.

The present model maintains that lawful physical conditions are actual independently of whether a particular observer has received information about them.

This does not introduce instantaneous information transfer. Bob cannot use the universal-present proposition to learn the remote condition of the star or communicate with it faster than \(c\).

His observable information remains constrained by the causal structure of spacetime.

The earlier Timeless Light Model interpretation of \(c\) as a registration bound \cite{cBound} emphasizes the role of the invariant speed limit in governing physical information exchange. The present argument applies the same distinction between causal registration and physical actuality.

A limitation on what an observer can know is not, merely by being a limitation on knowledge, a limitation on what physically exists.

The universe is not required to await Bob's observation before possessing a lawful physical condition.

The information speed limit governs the accessibility and propagation of information. It does not serve as a universal clock determining which parts of the universe are permitted to be actual.

\section{Universal Actuality and Relativistic Simultaneity}

The proposal that there is one universal now must be distinguished from a conventional assignment of simultaneity to distant events.

In special relativity, the time coordinate assigned to a distant event depends on the observer's reference frame.

For a Lorentz transformation between inertial coordinate systems,

\begin{equation}
t'=\gamma\left(t-\frac{vx}{c^2}\right),
\end{equation}

where

\begin{equation}
\gamma=\frac{1}{\sqrt{1-v^2/c^2}}.
\end{equation}

Two spatially separated events assigned the same coordinate time in one inertial frame generally receive different time coordinates in another frame moving relative to the first.

This is the familiar relativity of simultaneity.

It is a genuine feature of relativistic spacetime description, not simply a failure to correct for light-travel delay.

The present paper does not deny that feature.

Instead, it distinguishes coordinate simultaneity from the ontological claim that physical actuality is singular.

Coordinate simultaneity is a relationship assigned to events within a specified spacetime description. The universal present proposed here is not defined by assigning every distant physical condition the same numerical coordinate time.

Nor is it defined by requiring all clocks to agree, selecting a preferred inertial observer, or imposing a universal synchronization procedure.

It is the claim that physical existence is not divided into independently existing past, present, and future universes, or into different existing presents belonging to different observers.

The absence of a universal simultaneity coordinate does not logically establish the existence of multiple physical actualities.

Likewise, the existence of only one actuality does not, by itself, supply a unique relativistic simultaneity relation between distant events.

This distinction is essential to the proposed ontology.

When the paper says that the universe has one current physical configuration, it does not mean that the configuration has been supplied with a preferred global coordinate-time label. It means that physical actuality is singular and relational, rather than a sequence of completed physical universes simultaneously retained in existence.

The universal-now proposition is therefore an ontological premise. It is not presented as a theorem derived from Lorentz transformations, nor as a new experimentally established global foliation of spacetime.

A further theory that assigned a physically unique simultaneity relation to every distant pair of events would be making an additional claim. The present model makes no such coordinate assignment.

It requires only that differences in clock accrual, signal propagation, and observer-dependent simultaneity descriptions not be mistaken for proof of separate existing physical presents.

Under this distinction, all physical actuality belongs to the one present, even though no observer possesses a master clock or immediate access to all distant physical conditions.

\section{The Absence of a Master Spatial Grid}

The same issue arises in spatial description.

A coordinate system is ordinarily constructed by selecting an origin, directions, and rules for assigning coordinates to events or physical systems.

Such a system is useful because lawful relationships can be represented mathematically.

An observer may choose one star as an origin and construct a map of surrounding stars. The resulting map permits calculations of separations, relative velocities, and predicted interactions.

Another observer may choose a different origin and produce a different map of the same physical relationships.

Neither map thereby becomes the independent spatial reality underlying the stars.

The coordinate assignments are descriptions of relationships, not additional entities existing alongside those relationships.

The proposed universal present does not restore the independently existing spatial container rejected earlier in the paper.

A universe-wide current physical configuration is not a collection of independent substances placed at absolute coordinates on a master grid. It is the lawful relational organization of physical actuality.

The completeness of that actuality does not imply that one observer possesses a complete representation of it.

Indeed, the invariant information speed limit prevents any local observer from obtaining instantaneous reports from arbitrarily distant physical systems.

The configuration is therefore not identified with a completed observational map.

It is the physical actuality to which different local observational maps refer.

This gives spatial coordinates and local clock readings parallel roles. Both provide lawful quantitative descriptions. Neither is an independently existing master structure governing physical existence.

The absence of a master spatial grid does not mean the absence of spatial relationships.

The absence of a master clock does not mean the absence of locally accumulated proper time.

And the absence of a uniquely assigned distant simultaneity coordinate does not, within the proposed ontology, require the absence of one physical actuality.

\section{Worldlines as Descriptions of Lawful History}

The familiar representation of motion uses a worldline.

A timelike worldline represents the succession of spacetime events attributed to a physical system. Its geometry determines the proper-time interval accumulated along the represented history.

The present interpretation retains that mathematical representation.

It does not, however, require the entire represented worldline to exist as a completed physical trail.

A diagram may show Bob departing from one event, experiencing a series of interactions, and arriving at another event. The mathematical description organizes the lawful history associated with those events.

Its success as a description does not establish that earlier portions of the worldline remain physically in existence after the corresponding events.

The distinction is analogous to the distinction between a map and the physical relationships represented by the map.

A spatial map does not establish an independently existing spatial container. A worldline diagram does not, merely by representing an ordered history, establish an independently existing collection of past physical configurations.

The physical evidence supporting a reconstructed worldline consists of present records, present physical states, and the lawful relationships used to infer or calculate the represented history.

This distinction does not invalidate the calculation of proper-time intervals. The geometric quantity associated with a worldline remains mathematically defined, and clock comparisons remain physically meaningful.

It instead addresses the ontological interpretation of the representation.

The worldline is a lawful mathematical description of physical history. It need not be regarded as a permanently existing physical road through which a substance travels.

The same distinction applies to the worldlines of individual constituents within an organized physical system.

Their differing accumulated intervals describe differing histories. Those histories do not require the earlier configurations of the constituents to remain physically present as completed objects.

The records and physical consequences of those histories participate in current actuality.

The prior events themselves are not required to persist as independent physical entities.

\section{A Unified Relational Account of Space, Motion, and Time}

The preceding arguments converge on a common interpretation.

Space is not an independently existing container. What are ordinarily called objects are stable lawful physical organizations, and their positions are relational assignments.

Motion is the lawful evolution of spatial relationships, not independent substance displacement.

Proper time is accumulated along individual worldlines. The accumulated intervals depend on worldline geometry and the spacetime metric, rather than on a universal master clock.

The universe has one current physical actuality. Past events are represented through present records and lawful reconstructions, not through the continued physical existence of earlier completed universes.

The information speed limit governs the propagation of signals and the availability of information about distant conditions. It does not convert limitations on observation into multiple independently existing physical presents.

These statements describe different aspects of the same proposed ontology.

The spatial description explains lawful relationships among physical systems. The temporal description explains the accumulated proper-time histories associated with those systems. The causal description explains the lawful restrictions governing information propagation and physical interaction.

None requires a separately existing space or time container.

Nor does the account require the universe to synchronize the accrued histories of its constituents before allowing them to interact.

When Bob and Alice meet, their physical encounter is governed by their lawful local conditions. Their different accumulated histories remain real and measurable, but those differences do not constitute a conflict in physical actuality.

Likewise, when an observer receives a signal from a distant star, the registered signal has its lawful causal history. The delay does not establish that the star exists in a separate physical present.

The universe is described through one physically actual relational condition, local histories of proper-time accrual, and lawful changes and interactions.

The account thus replaces the image of objects traveling through space while time passes universally with an account of physical actuality as lawful relational organization.

The replacement is interpretive, not mathematical.

The equations of relativity remain intact because they already provide the quantitative description of the relationships, worldlines, and local accruals at issue.

\section{Scope and Logical Status of the No-Go}

The argument advanced here has two distinguishable foundations.

The first consists of established relativistic structure. Spatial velocity is frame-dependent. Proper acceleration is locally measurable. Proper time is associated with timelike worldlines. Different worldlines may accumulate different proper-time intervals between common events. Coordinate assignments, including assignments of distant simultaneity, depend on the chosen description.

The second foundation is ontological.

The paper adopts the premise that physical actuality consists of lawful relationships and configurations rather than independent substances occupying an independently completed spatial container. It further adopts the premise that there is one present physical actuality, without a separately existing past or future.

These premises are not established merely by the coordinate properties of special or general relativity.

They supply the interpretation under examination.

The no-go on independent spatial displacement follows conditionally from the relational premise. If there are no independently existing spatial positions or traveler-substances, there is no additional independent substance displacement for the motion description to identify.

The universal-present interpretation likewise follows from the adopted account of physical actuality. It does not establish a preferred relativistic reference frame, a new synchronization law, or an empirically distinguishable absolute simultaneity coordinate.

The paper therefore does not claim to have proved that alternative ontologies are mathematically impossible.

It argues that the independent-container and traveler-substance interpretations are unnecessary to the physical description, and that a singular-present, locally accrued, relational ontology provides a coherent alternative.

No prediction of a trajectory, orbital motion, Doppler shift, gravitational redshift, clock comparison, or local acceleration measurement is changed by adopting the interpretation.

Its intended contribution is a clarified account of what the successful mathematical descriptions require physically, and what additional entities they do not require.

\section{Conclusion}

Motion does not require independently existing objects traveling through an independently existing space.

Physical systems may instead be understood as stable patterns of lawful organization, with positions and velocities expressing relationships among those systems.

Bob is stationary in Bob's own coordinates. Alice is stationary in Alice's own coordinates. Their changing separation is a real physical relationship, and the lawful consequences of that change are registered through interaction.

Neither observer's description identifies an absolute position or intrinsic spatial velocity.

The conventional picture of a substance traversing a completed spatial container is therefore unnecessary under the relational ontology.

The same physical account extends to time.

Every physical constituent accumulates its own proper-time history along its worldline. Those histories are governed by the spacetime metric, including the effects associated with relative motion and gravitational conditions. Different histories do not require different physically existing presents.

There is one physical actuality, which we call now.

The past does not remain as a separately existing physical universe. Records and memories exist within present actuality, while the events they describe are reconstructed through lawful physical relationships.

The invariant information speed limit governs interactions and the accessibility of distant information. It does not establish that distant physical conditions lack actuality merely because their current states are inaccessible to a particular observer.

The resulting ontology requires neither a master spatial grid nor a universal elapsed-time controller.

The central result is therefore not that motion, time, or physical change is unreal. Their measurable consequences are real and lawful.

What is unnecessary is the additional picture of independent substances traveling through a completed space while a universal temporal substance passes around them.

There is physical actuality, consisting of lawful relationships and configurations. Those relationships change. Physical constituents accumulate their own proper-time histories. Observers register the consequences through lawful interactions.

Motion is relational. Proper-time accrual is local. Physical actuality is singular.

\section*{Relation to the Timeless Light Model}

The Timeless Light Model distinguishes lawful physical relationships from the additional traveler-substance interpretation ordinarily attached to physical propagation. Its photon-side interpretation treats light as a lawful relation expressed through endpoint registrations rather than requiring a particle-substance to traverse an independently occupied middle.

The present paper extends the same ontological discipline to material systems and their spatial relationships.

A physical system need not be understood as a substance traveling through an independently completed spatial interval. Its observable motion-content is the lawful change of relationships among physical systems.

The local accumulation of proper time remains associated with timelike worldlines, while causal registration remains constrained by the invariant information speed limit.

The resulting account places spatial motion, local accrual, and the universal present within one relational interpretation of physical actuality.

The traveler-substance and independent spatial container are unnecessary. The lawful relationships and their physical consequences remain.

\begin{thebibliography}{99}
\bibitem{SpaceContainer}
J.~C.~W.~McKinley.
\newblock \emph{Space Was Never a Completed Container}.
\newblock Working manuscript (2026).

\bibitem{Atom}
J.~C.~W.~McKinley.
\newblock \emph{The Atom Is Available When Called On: Stationary States as the Matter-Side Instance of General Timelessness}.
\newblock Zenodo (2026).
\newblock \doi{10.5281/zenodo.20114822}.

\bibitem{Accrual}
J.~C.~W.~McKinley.
\newblock \emph{Local Accrual as the Only Intrinsic Time-Quantity: A Positive Structural Statement for the Massive Regime}.
\newblock Zenodo (2026).
\newblock \doi{10.5281/zenodo.20225645}.

\bibitem{DilationProperty}
J.~C.~W.~McKinley.
\newblock \emph{Dilation Is Not a Property of the Object: A No-Go Result on Intrinsic Time-Dilation}.
\newblock Zenodo (2026).
\newblock \doi{10.5281/zenodo.20225720}.

\bibitem{WorldlinesProduce}
J.~C.~W.~McKinley.
\newblock \emph{Worldlines Produce; Dilation Reports: A No-Go on Treating Consequence-Reality as Evidence of Intrinsic Dilation}.
\newblock Zenodo (2026).
\newblock \doi{10.5281/zenodo.20225757}.

\bibitem{cBound}
J.~C.~W.~McKinley.
\newblock \emph{c Is a Registration Bound, Not a Traveler's Speed: A Registration-Based Interpretation of the Information Speed Limit}.
\newblock Zenodo (2026).
\newblock \doi{10.5281/zenodo.20114176}.

\bibitem{NoPast}
Ellison Clarke.
\newblock \emph{The Universe Has No Past}.
\newblock Working manuscript (2026).

\bibitem{Einstein1905}
A.~Einstein.
\newblock Zur Elektrodynamik bewegter K{\"o}rper.
\newblock \emph{Annalen der Physik}, \textbf{17}, 891--921 (1905).

\bibitem{Wald}
R.~M.~Wald.
\newblock \emph{General Relativity}.
\newblock University of Chicago Press (1984).

\bibitem{TaylorWheeler}
E.~F.~Taylor and J.~A.~Wheeler.
\newblock \emph{Spacetime Physics}.
\newblock W.~H.~Freeman, 2nd edition (1992).
\end{thebibliography}
\end{document}
